% ============================================================
%  cap06.tex  --  Capítulo 6: Formas de Cotar a Taxa de Juros
%  Baseado em: Notas de Aula 7 (Cap. 7)
%  Referência: Bueno, Rangel & Santos, Cap. 7
% ============================================================

\chapter{Formas de Cotar a Taxa de Juros}

\section{Objetivo}

Este capítulo tem dois objetivos complementares:
\begin{enumerate}[label=\textbf{\arabic*.}]
  \item Distinguir \textbf{taxa nominal} de \textbf{taxa efetiva}.
  \item Introduzir o conceito de \textbf{taxa real} (descontada a inflação).
\end{enumerate}

Confundir essas cotações é fonte frequente de erro em contratos, projeções financeiras e análise de investimentos.

% ─────────────────────────────────────────────────────────────────────────────
\section{Convenções de Contagem de Dias}

\begin{center}
\renewcommand{\arraystretch}{1.3}
\begin{tabular}{lll}
\toprule
\textbf{Convenção} & \textbf{Ano-base} & \textbf{Mês-base} \\
\midrule
Dias corridos & 360 dias & 30 dias \\
Dias úteis    & 252 dias úteis & 21 dias úteis \\
\bottomrule
\end{tabular}
\end{center}

\begin{nota}
Na prática, devido a feriados, nem todos os meses têm o mesmo número de dias úteis.
A estratégia correta é sempre utilizar os dias \emph{exatos} do período.

\smallskip
\textbf{Ferramentas:}
\begin{itemize}[leftmargin=1.5em, nosep]
  \item \textbf{HP 12C}: calcula dias efetivos entre duas datas (\texttt{g\;$\Delta$DYS}),
        base 360 ou 365 dias.
  \item \textbf{Excel}: \texttt{DIAS360} (corridos) ou \texttt{DIATRABALHOTOTAL}
        (úteis, requer tabela de feriados inserida pelo usuário).
\end{itemize}
\end{nota}

% ─────────────────────────────────────────────────────────────────────────────
\section{Taxa Nominal vs.\ Taxa Efetiva}

\begin{definicao}
A \textbf{taxa nominal} é uma taxa anunciada para um período de referência que \emph{não}
corresponde diretamente à capitalização efetiva. Por convenção, ela é proporcional à
subperiodização.

A \textbf{taxa efetiva} é a taxa que de fato remunera o capital levando em conta a capitalização
real no período analisado.
\end{definicao}

Se a taxa nominal anual é $j$ com capitalização $m$ vezes por ano, a taxa efetiva anual é:
\[
  i_{\text{ef}} = \left(1 + \frac{j}{m}\right)^m - 1.
\]

\begin{exemplo}{Selic e taxa efetiva (exemplo brasileiro)}
A taxa Selic-meta é anunciada como uma taxa \emph{anual} mas capitaliza diariamente
(base 252 dias úteis). Uma taxa anual de $6\%$ com capitalização diária em base 252:
\[
  i_m = (1 + 0{,}06/252)^{21} - 1 \approx 0{,}502\% \text{ a.m.}
\]
Já a taxa de $0{,}5\%$ a.m.\ anunciada como nominal (taxa proporcional):
\[
  i_{\text{ef}} = (1{,}005)^{12} - 1 \approx 6{,}17\% \text{ a.a.}
\]
\end{exemplo}

\subsection{Equivalência de taxas}

O princípio de equivalência (visto no Cap.~1) aplica-se diretamente:
\[
  (1 + i_1)^{n_1/360} = (1 + i_2)^{n_2/360}
\]
Para converter entre dias corridos e dias úteis, usa-se o número real de dias do período.

% ─────────────────────────────────────────────────────────────────────────────
\section{Taxa de Juros Real}

\subsection{Motivação}

Até aqui ignoramos a inflação (variação geral do nível de preços). Uma taxa de juros nominal
de $12\%$ com inflação de $8\%$ não rende $12\%$ em termos de poder de compra real.

\subsection{Derivação da Equação de Fisher}

Considere um ativo de preço $P_0$ hoje e $P_1 = P_0(1+\pi)$ no próximo período
(onde $\pi$ é a inflação). Um investimento $V_0$ cresce para $V_1 = V_0(1+i)$ em termos nominais.
O poder de compra real ao final é:
\[
  \frac{V_1}{P_1} = \frac{V_0(1+i)}{P_0(1+\pi)}
                  = \frac{V_0}{P_0} \cdot \frac{1+i}{1+\pi}.
\]
Definindo a taxa de retorno real $r$ como:
\[
  1 + r = \frac{1+i}{1+\pi}
\]

\begin{resultado}
\[
  \boxed{1 + r = \frac{1+i}{1+\pi}}
  \qquad\Leftrightarrow\qquad
  r = \frac{i - \pi}{1 + \pi}.
\]
Esta é a \textbf{Equação de Fisher}.
\end{resultado}

\paragraph{Aproximação (para $i$ e $\pi$ pequenos):}
\[
  r \approx i - \pi.
\]

\begin{exemplo}{Taxa real -- referência brasileira 2024}
Taxa Selic: $i = 10{,}5\%$ a.a. \quad IPCA: $\pi = 4{,}62\%$ a.a.
\[
  r = \frac{1{,}105 - 1{,}0462}{1{,}0462} - 1
    = \frac{1{,}105}{1{,}0462} - 1
    \approx 5{,}62\% \text{ a.a. (real).}
\]
Pela aproximação: $r \approx 10{,}5\% - 4{,}62\% = 5{,}88\%$ (ligeiramente acima do valor exato).
\end{exemplo}

\subsection{Taxa real com capitalização contínua}

No regime contínuo, com taxas nominais $r_{\text{nom}}$ e inflação $\pi_c$ (ambas contínuas):
\[
  r_{\text{real}} = r_{\text{nom}} - \pi_c.
\]
A subtração simples é exata em tempo contínuo -- não há o termo $r\pi$ de correção.

\subsection{Caso mais realista: inflação variável}

Se a inflação $\pi_t$ varia a cada período $t = 1, \ldots, n$:
\[
  (1+r)^n = \frac{(1+i)^n}{\prod_{t=1}^n(1+\pi_t)}.
\]
A taxa real ``média geométrica'' é:
\[
  1 + r = \frac{(1+i)}{\left[\prod_{t=1}^n(1+\pi_t)\right]^{1/n}}.
\]

% ─────────────────────────────────────────────────────────────────────────────
\section{Fluxo de Caixa Real}

Ao avaliar um projeto com fluxos em termos \textbf{nominais} $R_t$ usando a TIR nominal,
os resultados podem ser enganosos se não considerarmos a inflação. A alternativa é
trabalhar em \textbf{termos reais}:

\[
  R_t^{\text{real}} = \frac{R_t}{(1+\pi_1)(1+\pi_2)\cdots(1+\pi_t)}
                     \approx \frac{R_t}{(1+\bar\pi)^t},
\]

e descontar pela taxa real $r$:

\[
  P = \sum_{t=1}^n \frac{R_t^{\text{real}}}{(1+r)^t}.
\]

\begin{nota}
Em geral usamos o IPCA/FGV como índice de reajuste $\pi_t$.
O resultado é o mesmo se descontarmos o fluxo nominal pela taxa nominal $i$,
\emph{desde que} a taxa $i$ e os fluxos $R_t$ sejam consistentes (ambos nominais ou ambos reais).
\end{nota}

\begin{center}
\begin{tikzpicture}[>=Stealth, font=\small, scale=0.95]
  \draw[timeline axis] (-0.5,0)--(8.5,0) node[right]{$t$};
  % Ticks
  \foreach \x/\lab in {0/0,1/1,2/2,3/3,7/$n$}{
    \draw[tick] (\x,0.12)--(\x,-0.12);
    \node[below=4pt] at (\x,0) {\lab};
  }
  \node at (5.5,-0.35) {$\cdots$};
  % P/p_0 (upward)
  \draw[arrow up] (0,0)--(0,1.4);
    \node[above=2pt,color=red!70!black] at (0,1.4) {$P_0/p_0$};
  % Fluxos reais R_t/p_t (downward)
  \draw[arrow down] (1,-0.65)--(1,-1.4);
    \node[cash label,below=2pt] at (1,-1.45) {$\frac{R_1}{p_1}$};
  \draw[arrow down] (2,-0.65)--(2,-1.4);
    \node[cash label,below=2pt] at (2,-1.45) {$\frac{R_2}{p_2}$};
  \draw[arrow down] (3,-0.65)--(3,-1.4);
    \node[cash label,below=2pt] at (3,-1.45) {$\frac{R_3}{p_3}$};
  \draw[arrow down] (7,-0.65)--(7,-1.4);
    \node[cash label,below=2pt] at (7,-1.45) {$\frac{R_n}{p_n}$};
\end{tikzpicture}
\end{center}

% ─────────────────────────────────────────────────────────────────────────────
\section{Perpetuidade com Inflação (Fisher + Gordon)}

Para uma perpetuidade cujos pagamentos crescem pela inflação $\pi$ (série geométrica
com $\alpha = \pi$):
\[
  R_1^{\text{nominal}} = R_1^{\text{real}} \cdot (1+\pi).
\]
Usando a fórmula de Gordon:
\[
  P = \frac{R_1}{i - \pi}.
\]
Pela equação de Fisher, $i - \pi \approx r(1+\pi)$, logo:
\[
  P \approx \frac{R_1}{r(1+\pi)} = \frac{R_1^{\text{real}}}{r}.
\]

\begin{resultado}
A perpetuidade nominal crescente à inflação $\pi$ vale o mesmo que a perpetuidade
real uniforme com fluxo $R^{\text{real}}$ descontada pela taxa real $r$.
\[
  P = \frac{R_1}{i-\pi} = \frac{R^{\text{real}}}{r}.
\]
\end{resultado}

\newpage

\begin{moderno}[title={Tesouro IPCA+ como perpetuidade real (2025)}]
O título NTN-B paga cupom semestral de $6\%$ a.a.\ mais correção pelo IPCA.
Para uma NTN-B de prazo muito longo (``NTN-B Perpétua'' -- hipotética), com
taxa real demandada $r = 6\%$ a.a.\ e cupom semestral $c = 3\%$ ($\approx 6\%/2$) sobre
o valor de face de R\$\,1.000 (corrigido pelo IPCA):
\[
  P \approx \frac{30}{0{,}03} = \text{R\$\,}1{.}000
  \qquad\text{(à par, quando a taxa de mercado = cupom).}
\]
Se o mercado exige $r = 7\%$ a.a., o preço cai para $\approx$ R\$\,857 (abaixo do par).
\end{moderno}
